\scanpage{35}\mspage{41}
% archive sheet number: 34
\deleted{Donc si
\[
 \operatorname{hom}(e,F)\longrightarrow\operatorname{hom}(X,F)
\]
est iso pour tout $F$ dans $\mathrm{Hot}$, \unsure{passant|prenant} aux $\pi_0$ on en déduit que
\[
 \operatorname{Hom}_{\mathrm{Hot}}(e,F)\longrightarrow\operatorname{Hom}_{\mathrm{Hot}}(X,F)
\]
est iso pour \unsure{tout} $F$, donc $X\to e$ iso dans $\mathrm{Hot}$, i.e. $X$ contractile, OK.

Mais le cas général \unsure{de la|du} ``\unsure{propr.|principe}'' reste douteux \ldots}

\margin{\uline{NB} les pages 33 à 40 ont été supprimées, car \unsure{canulées}}

Je vais reprendre la terminologie et les axiomes relatifs aux dérivateurs.

Un \uline{prédérivateur} est un 2-foncteur
\[
 \mathrm{Diag}^{\circ}\xrightarrow{\mathbb D}(\mathrm{Cat}),
\]
% slanted
\margin{Diag : \uline{domaine} du dérivateur. Les objets sont les \uline{cat. d'}\deleted{\unsure{coefficients}}\inserted{indices} pour $\mathbb D$. Les $\mathbb D(I)$ sont les cat. de \unsure{coeffs} pour $\mathbb D$, et les objets de $\mathbb D(I)$ \unsure{des} ``\unsure{coefficients|coeffts} sur $I$''}
$(\mathrm{Diag})$ est une \inserted{2-}sous-catégorie pleine de $(\mathrm{Cat})$. On \unsure{la supposera|suppose} tout au moins que Diag contient les \uline{ens.} \uline{ordonnés} \uline{finis}, que Diag est stable par $\underleftarrow{\lim}$ finies, par sommes finies. Qu'avec toute catégorie $C$ et toute sous-catégorie strictemt pleine ouverte (resp. fermée) elle contienne la sous-catégorie complémentaire.
